\let\negmedspace\undefined
\let\negthickspace\undefined
\documentclass[journal]{IEEEtran}
\usepackage[a5paper, margin=10mm, onecolumn]{geometry}
%\usepackage{lmodern} % Ensure lmodern is loaded for pdflatex
\usepackage{tfrupee} % Include tfrupee package

\setlength{\headheight}{1cm} % Set the height of the header box
\setlength{\headsep}{0mm}     % Set the distance between the header box and the top of the text

\usepackage{gvv-book}
\usepackage{gvv}
\usepackage{cite}
\usepackage{amsmath,amssymb,amsfonts,amsthm}
\usepackage{algorithmic}
\usepackage{graphicx}
\usepackage{textcomp}
\usepackage{xcolor}
\usepackage{txfonts}
\usepackage{listings}
\usepackage{enumitem}
\usepackage{mathtools}
\usepackage{gensymb}
\usepackage{comment}
\usepackage[breaklinks=true]{hyperref}
\usepackage{tkz-euclide} 
\usepackage{listings}
% \usepackage{gvv}                                        
\def\inputGnumericTable{}                                 
\usepackage[latin1]{inputenc}                                
\usepackage{color}                                            
\usepackage{array}                                            
\usepackage{longtable}                                       
\usepackage{calc}                                             
\usepackage{multirow}                                         
\usepackage{hhline}                                           
\usepackage{ifthen}                                           
\usepackage{lscape}
\usepackage{circuitikz}



\author{EE25BTECH11041-Naman Kumar }
\graphicspath{./figs/}

\begin{document}
\begin{center}
    \huge{4.3.30}\\
    \large{EE25BTECH11041 - Naman Kumar}
\end{center}
Question:\\
Find the equation of the line which passes through the point $(-4, 3)$ and the portion of the line intercepted between the axes is divided internally in ratio 5:3 by this point.\\
\solution \\
Let the intercept points be
\begin{align}
    \vec{P}=\begin{pmatrix}a\\0\end{pmatrix},\vec{Q}=\begin{pmatrix}0\\b\end{pmatrix}\text{ and }\vec{R}=\begin{pmatrix}-4\\3\end{pmatrix}
\end{align}
be the given point.\\
Using
\begin{align}
    \vec{R}=\frac{k\vec{Q}+\vec{P}}{k+1}
\end{align}
where k=$\frac{5}{3}$\\
Using Putting all the values
\begin{align}
    \vec{R}=\frac{5\times \begin{pmatrix}0\\b\end{pmatrix}+3\times \begin{pmatrix}a\\0\end{pmatrix}}{8}\\
    8\times\begin{pmatrix}-4\\3\end{pmatrix}=\begin{pmatrix}3a\\5b\end{pmatrix}\\
    \begin{pmatrix}-32\\24\end{pmatrix}=\begin{pmatrix}3a\\5b\end{pmatrix}\\
    \begin{pmatrix}a\\b\end{pmatrix}=\begin{pmatrix}\frac{-32}{3}\\ \frac{24}{5}\end{pmatrix}
\end{align}
therefore, Intercepts of line are
\begin{align}
    \vec{P}=\begin{pmatrix}\frac{-32}{3}\\0\end{pmatrix},\vec{Q}=\begin{pmatrix}0\\ \frac{24}{5}\end{pmatrix}
\end{align}
General equation of line
\begin{align}
    \vec{x}=\vec{h}+k\vec{m}
\end{align}
Where\\
\input{tables/table.tex}
Slope is
\begin{align}
\vec{m}=\vec{Q}-\vec{P}\\
\vec{m}=\begin{pmatrix}\frac{32}{3}\\ \frac{24}{5}\end{pmatrix}
\end{align}
let $\vec{h}=\vec{P}$\\
So, Equation of line is
\begin{align}
    \vec{x}=\vec{h}+k\vec{m}\\
    \begin{pmatrix}x\\y\end{pmatrix}=\begin{pmatrix}\frac{-32}{3}\\0\end{pmatrix}+k\begin{pmatrix}\frac{32}{3}\\ \frac{24}{5}\end{pmatrix}
\end{align}
\begin{figure}[H]
    \centering
    \includegraphics[width=\columnwidth]{figs/figure.png}
    \label{fig:placeholder}
\end{figure}
\end{document}
